% Part: first-order-logic
% Chapter: proof-systems
% Section: seqeunt-calculus

\documentclass[../../../include/open-logic-section]{subfiles}

\begin{document}

\iftag{FOL}
      {\olfileid{fol}{prf}{seq}}
      {\olfileid{pl}{prf}{seq}}

\olsection{સિક્વન્ટ કલન}

ઘણાં !!{derivation} તંત્રો !!{sentence}sની ગોઠવણી પર કાર્ય કરે છે,
જ્યારે સિક્વન્ટ કલન \emph{સિક્વન્ટો} પર કાર્ય કરે છે. સિક્વન્ટ નીચેના
સ્વરૂપનું પદ છે:
\[
!A_1, \dots, !A_m \Sequent !B_1, \dots, !B_n,
\]
એટલે કે, સિક્વન્ટ-સંકેત~$\Sequent$થી અલગ કરેલી !!{sentence}sની બે
શ્રેણીઓની જોડી. બંનેમાંથી કોઈપણ શ્રેણી ખાલી હોઈ શકે છે. સિક્વન્ટ કલનમાં
!!^a{derivation} સિક્વન્ટોનું વૃક્ષ હોય છે. તેમાં સૌથી ઉપરના સિક્વન્ટો
ખાસ સ્વરૂપના હોય છે (તેમને ``આરંભિક સિક્વન્ટો'' અથવા
``સ્વયંસિદ્ધિઓ'' કહે છે), અને દરેક બીજો સિક્વન્ટ તેની બરાબર ઉપરના
સિક્વન્ટોમાંથી અનુમાનના કોઈ એક નિયમ વડે મળે છે. અનુમાનના નિયમો કાં તો
સિક્વન્ટોમાંનાં !!{sentence}sમાં ફેરફાર કરે છે (તેમને ડાબી કે જમણી
બાજુ ઉમેરે, દૂર કરે અથવા ફરી ગોઠવે છે), અથવા નિયમના ફલિતમાં કોઈ
સંકુલ !!{formula} દાખલ કરે છે. દાખલા તરીકે, નિયમ $\LeftR{\land}$ વડે
$!A, \Gamma \Sequent \Delta$ પરથી $!A \land !B, \Gamma \Sequent
\Delta$નું અનુમાન કરી શકાય છે; અને નિયમ $\RightR{\lif}$ વડે કોઈપણ
$\Gamma$, $\Delta$, $!A$ અને~$!B$ માટે $!A, \Gamma \Sequent \Delta,
!B$ પરથી $\Gamma \Sequent \Delta, !A \lif !B$નું અનુમાન કરી શકાય છે.
(ખાસ કરીને, $\Gamma$ અને~$\Delta$ ખાલી હોઈ શકે છે.)
\sourcecorrection{OLPRF-001}{મૂળની \texttt{sequent-calculus.tex},
પંક્તિ~19માં સિક્વન્ટની સ્વતંત્ર ડાબી અને જમણી શ્રેણી બંનેનો છેલ્લો
ક્રમસૂચક $m$ લખાયો છે. આ જ ગ્રંથની સિક્વન્ટની વિસ્તૃત વ્યાખ્યા મુજબ
ડાબે $!A_1,\dots,!A_m$ અને જમણે $!B_1,\dots,!B_n$ લખીને બંને
શ્રેણીઓની લંબાઈ સ્વતંત્ર રાખી છે.}

સિક્વન્ટ કલન પર આધારિત $\Proves$ સંબંધ આ પ્રમાણે વ્યાખ્યાયિત થાય છે:
$\Gamma \Proves !A$ જો અને માત્ર જો એવી કોઈ શ્રેણી $\Gamma_0$ હોય કે
$\Gamma_0$માંનું દરેક $!A$ એ~$\Gamma$માં હોય અને જેના મૂળે
$\Gamma_0 \Sequent !A$ સિક્વન્ટ હોય એવી !!{derivation} હોય. જો
$\Sequent !A$ સિક્વન્ટને !!a{derivation} હોય, તો સિક્વન્ટ કલનમાં $!A$
પ્રમેય છે. દાખલા તરીકે, નીચેની !!a{derivation} બતાવે છે કે
$\Proves (!A \land !B) \lif !A$:
\begin{prooftree}
  \Axiom$!A \fCenter !A$
  \RightLabel{\LeftR{\land}}
  \UnaryInf$!A \land !B \fCenter !A$
  \RightLabel{\RightR{\lif}}
  \UnaryInf$\fCenter (!A \land !B) \lif !A$
\end{prooftree}

જો $\Gamma_0 \Sequent$ની !!a{derivation} હોય (જ્યાં દરેક
$!A \in \Gamma_0$ એ~$\Gamma$માં હોય અને સિક્વન્ટની જમણી બાજુ ખાલી
હોય), તો ગણ $\Gamma$ સિક્વન્ટ કલનમાં અસુસંગત છે. નિયમ
\RightR{\Weakening} વડે અસુસંગત ગણમાંથી કોઈપણ !!{sentence}
!!{derive}d કરી શકાય છે.

સિક્વન્ટ કલનની શોધ ગેરહાર્ડ ગેન્ટ્ઝને 1930ના દાયકામાં કરી હતી. તેની
પદ્ધતિસર અને સમમિત રચનાને કારણે !!{derivation}sનો સિદ્ધાંત વિકસાવવા
માટે તે બહુ ઉપયોગી ઔપચારિક પદ્ધતિ છે. સિક્વન્ટ કલનમાં
!!{derivation}s શોધવી પ્રમાણમાં સહેલી છે, પરંતુ આ !!{derivation}s
વાંચવી ઘણી વાર અઘરી હોય છે અને સાબિતીઓ સાથેનો તેમનો સંબંધ ક્યારેક
સરળતાથી દેખાતો નથી. તેમ છતાં, !!{derivation} તંત્રો માટે તે અત્યંત
સુઘડ અભિગમ સાબિત થયો છે, અને ઘણાં તર્કશાસ્ત્રોને સિક્વન્ટ-કલન તંત્રો
છે.

\end{document}
